\begin{thebibliography}{99}
\markboth{参考図書}{参考図書}
\addcontentsline{toc}{chapter}{参考図書}

\bibitem{ryden2025}
Barbara Ryden, ``\textit{Celestial and Stellar Dynamics}'', Cambridge University Press, 2025.\\
学部生向けに天体力学と恒星系力学を扱う入門的な教科書。
本教材では、物理的な題材の主な出典として、§6.4--6.5 の球対称・軸対称系と、§7.1--7.3 の球対称・軸対称・非軸対称ポテンシャル中の軌道を参照する。

\bibitem{bt2008}
James Binney \& Scott Tremaine, ``\textit{Galactic Dynamics}'', 2nd ed., Princeton University Press, 2008.\\
銀河力学・恒星系力学を体系的に扱う標準的な教科書。
第2章の James 法の補足では、§2.9 の Box 2.5 を参照する。
付録Bの補足では、§2.4 の多重極展開、§2.8 の基底関数展開、§2.9 の $N$ 体計算における重力計算法を参照する。

\bibitem{james1977}
R. A. James, ``The solution of Poisson's equation for isolated source distributions'',
\textit{Journal of Computational Physics}, 25(2), 71--93, 1977.\\
有限の格子で孤立系の Poisson 方程式を解くため、
境界上の遮蔽質量と離散 Green 関数を利用した境界補正法を提案した論文。

\bibitem{fukue2022}
福江純・和田桂一・梅村雅之『宇宙流体力学の基礎［改訂版］』、日本評論社、2022年。\\
宇宙流体力学の物理的な基礎を体系的に扱う教科書。
宇宙流体の基礎方程式、自己重力流体、波動、衝撃波、流体不安定性、重力不安定性などを学ぶことができる。
本教材のあとに宇宙流体シミュレーションへ進む際に、数値計算で解く方程式とその物理的意味を学ぶための参考書とする。

\bibitem{kouya2021}
幸谷智紀『Python数値計算プログラミング』講談社、2021年。\\
Pythonによる数値計算法を、理論と実装の両面から扱う教科書。
本教材では、数値積分、数値微分、有限差分法、常微分方程式の数値解法を主に参照する。
付録Dでは、第7--8章の連立一次方程式、疎行列、反復法を参照する。
より入門的な説明には、杉浦『数値計算の基礎と応用―数値解析学への入門―』や、三井田・須田『数値計算法』も参考になる。

\bibitem{sugiura2009}
杉浦洋『数値計算の基礎と応用―数値解析学への入門―』新訂版、サイエンス社、2009年。\\
数値計算法の基礎を簡潔にまとめた入門書。
第1章では、数値積分則の次数、Taylor 展開による誤差評価、および局所的な近似誤差と大域的な誤差を区別する考え方の参考とする。

\bibitem{miida2023}
三井田惇郎・須田宇宙『数値計算法』第3版、森北出版、2023年。\\
数値計算法の基礎を具体的な計算例とともに学べる入門書。
第1章の補足では、モンテカルロ法と乱数生成、とくに逆変換法によるサンプリングの参考とする。

\bibitem{tomisaka2007}
富阪幸治・花輪知幸・牧野淳一郎 編『シミュレーション天文学』第2版、
シリーズ現代の天文学 第14巻、日本評論社、2026年。\\
常微分方程式・偏微分方程式の数値解法から、流体・磁気流体、自己重力流体、重力多体系まで、天文シミュレーションの基礎と応用を広く扱う教科書。
第2章では §5.3 の自己重力流体計算と多重極展開による境界値を参照する。
第1・3章の重力多体系に関する補足では第6章、付録Bの補足では重力計算法を参考とする。

\bibitem{matsumoto}
松本倫明・大須賀健・須佐元『輻射電磁流体シミュレーションの基礎』
シリーズ〈宇宙物理学の基礎〉第5巻、日本評論社、2024年。\\
宇宙流体シミュレーションの数値計算法を、格子法・粒子法から磁気流体・輻射流体まで広く扱う教科書。
第1章では §2.5 の有限体積法とセル平均、第2章では §2.4--2.6 の安定性・境界条件を参照する。
第3章では §9.6 のリープフロッグ法と長時間挙動を参照し、ソフトニング（§9.5）と流体の時間積分に関する補足でも参考とする。
付録Bでは、重力計算法に関する説明を参考とする。

\bibitem{fujii2020}
藤井孝藏・立川智章『Pythonで学ぶ流体力学の数値計算法』オーム社、2020年。\\
圧縮性流体の数値計算法を、Pythonによる実装とともに学べる教科書。
第2章末の偏微分方程式の型・境界条件・陽的および陰的時間積分の補足、第3章と Notebook 3 の RK2・RK3 の補足で参考とする。

\bibitem{aarseth2003}
Sverre J. Aarseth, ``\textit{Gravitational N-Body Simulations: Tools and Algorithms}'', Cambridge University Press, 2003.\\
重力 $N$ 体シミュレーションの数値計算法を詳しく扱う教科書。
直接重力計算、時間積分、近接遭遇、正則化、個別時間刻み、精度と計算効率などを学ぶことができる。
特に、星団のような衝突系の重力多体系計算を詳しく学ぶ際の参考になる。

\bibitem{dehnen2011}
Walter Dehnen \& Justin I. Read, ``\textit{N-body simulations of gravitational dynamics}'', \textit{European Physical Journal Plus}, 126, 55, 2011.\\
重力 $N$ 体シミュレーションの数値計算法を概説したレビュー。
衝突系と無衝突系の違いを踏まえ、それぞれで用いられる代表的な数値手法を整理している。
重力多体系計算の全体像を把握するための参考になる。

\end{thebibliography}
